Start from a finite set $A\subseteq\mathbb{N}$. Define a sequence $a_1<a_2<\cdots$ by taking $\{a_1,\dots,a_n\}$ to be the set of elements constructed so far (initially the given $A$), and then defining $a_{n+1}$ to be the smallest integer $>a_n$ which is \emph{not} of the form $a_i+a_j$ with $i,j\le n$ (repetitions allowed).
Equivalently, we repeatedly adjoin the smallest integer larger than the current maximum that does not lie in the sumset of the current set with itself.

The question: are the successive differences $a_{n+1}-a_n$ eventually periodic for every starting set $A$?
